Grandes modelos de linguagem (\textit{Large Language Models}, LLMs) deixaram de operar apenas como geradores de texto e passaram a compor agentes capazes de consultar dados, selecionar ferramentas e produzir chamadas que alteram sistemas. Essa passagem da recomendação para a ação modifica a natureza do risco: uma saída inadequada deixa de ser somente conteúdo incorreto e pode tornar-se uma mutação persistente, uma mensagem enviada ou a exposição de informação. Ambientes como AgentDojo e InjecAgent mostram que agentes conectados a ferramentas permanecem suscetíveis a instruções presentes em dados externos, ao mesmo tempo em que podem falhar em tarefas legítimas mesmo sem ataque \cite{debenedetti2024agentdojo,zhan2024injecagent}. Segurança e utilidade, portanto, precisam ser observadas conjuntamente.

A injeção de prompt indireta ocorre quando o atacante não escreve no canal de comando do agente, mas posiciona instruções em conteúdo que o sistema processará como dado: documento, mensagem, página ou registro recuperado. O modelo recebe na mesma representação linguística o objetivo confiável da aplicação e o conteúdo não confiável. Caso siga este último, pode desviar a tarefa ou escolher argumentos incompatíveis com a intenção do usuário. Greshake et al. demonstraram esse vetor em aplicações integradas a LLM e sistematizaram impactos que incluem chamadas indevidas de API e exfiltração \cite{greshake2023indirect}. Trabalhos posteriores construíram benchmarks mais amplos, mas também verificaram que defesas avaliadas apenas contra ataques estáticos podem ser contornadas por ataques adaptados ao mecanismo conhecido \cite{yi2025bipia,zhan2025adaptive}.

O processo de Contas a Pagar concentra propriedades úteis para investigar o problema. Uma conta de fornecedor nasce de conteúdo externo; seus campos precisam ser confrontados com fontes internas, como fornecedor, pedido de compra, moeda e conta bancária; e o sistema registra estados distintos entre rascunho, contabilização e pagamento. A própria documentação do Odoo separa a criação e o preenchimento de uma conta de fornecedor de sua confirmação, que gera o lançamento contábil \cite{odoo2026vendorbills}. Isso permite definir consequências mensuráveis sem movimentar dinheiro real: alterar campos financeiros, redirecionar o beneficiário indicado ou contabilizar sem autorização são eventos verificáveis por código.

O estudo procura produzir evidência contextualizada que relacione a chamada proposta pelo modelo ao efeito end-to-end em um processo empresarial. A comparação com trabalhos correlatos examina invariantes do ERP, separação entre rascunho e ação crítica, origem confiável dos campos e restauração do estado. Reduzir ataques bloqueando todas as ações também não constitui uma solução útil. A revisão dirigida e sua matriz crítica posicionam a contribuição como avaliação contextualizada e reproduzível desse recorte, sem reivindicar cobertura sistemática ou provar a inexistência de trabalhos equivalentes fora da seleção.

O problema de pesquisa é avaliar em que medida uma defesa em profundidade reduz violações no fluxo delimitado sem comprometer excessivamente as tarefas legítimas. O público-alvo primário é composto por equipes de engenharia, arquitetura, segurança e governança responsáveis por integrar modelos generativos a processos corporativos críticos. Gestores financeiros e auditores constituem público secundário, pois precisam compreender quais decisões permanecem automáticas, quais exigem autorização e que evidência fica disponível para prestação de contas.

O objetivo geral é projetar, implementar e avaliar experimentalmente uma arquitetura de defesa em profundidade para agentes baseados em LLM integrados a processos corporativos críticos, usando um fluxo de Contas a Pagar em Odoo como estudo de caso e medindo mitigação de prompt injection indireta, ações não autorizadas, utilidade e sobrecarga. A pergunta principal é:

\begin{quote}
\itshape Em que medida uma arquitetura de defesa em profundidade reduz o sucesso de ataques de prompt injection indireta e ações não autorizadas em um agente LLM integrado a um processo de Contas a Pagar, preservando a utilidade da automação e com qual custo operacional?
\end{quote}

Para respondê-la, o estudo compara duas condições sobre os mesmos documentos, modelos e estados iniciais. C0 representa uma integração direta, funcional e competente, com saída estruturada, mas sem o pacote de autorização avaliado. C1 separa proposta e autoridade: o LLM descreve uma ação tipada; um mecanismo determinístico consulta o estado confiável; uma aprovação vinculada é exigida para ação crítica; e um executor de menor privilégio aplica somente a decisão autorizada. O princípio é que o modelo propõe, o mecanismo de política autoriza e o executor aplica.

As contribuições pretendidas são: (i) uma arquitetura implementada que explicita fronteiras de confiança e autoridade; (ii) o conjunto AP-IPI-BR, formado por documentos sintéticos benignos e adversariais pareados; (iii) um protocolo que verifica o sucesso do ataque no estado do ERP, por observação independente do LLM; e (iv) evidência quantitativa de segurança e utilidade, acompanhada de descrição do custo e análise de falhas. A avaliação estima o efeito do pacote completo de proteção, sem isolar causalmente cada camada. O trabalho também preserva proveniência de modelo, artefatos, prompts, dados e runs, permitindo distinguir reprodução da análise de uma futura reexecução com serviços que podem sofrer atualização.

Esta versão registra o estágio atual do TCC II: além das provas técnicas do laboratório, dos caminhos C0/C1, da captura/restauração e do registro candidato das execuções, foi concluído o piloto de propostas em duas rodadas de 144 chamadas com o modelo real, nas quais nenhuma das 72 propostas adversariais de C0 incorporou o objetivo do atacante e as propostas benignas foram exatas em 36 de 36 execuções em C0 e 35 de 36 em C1. O ciclo integrado no estado do ERP para os pares do piloto, a decisão do grupo sobre o enquadramento estimativo, a revisão nominal, o congelamento e a coleta confirmatória permanecem pendentes. Essa distinção preserva a fronteira entre viabilidade técnica, propostas do piloto e efeito final no sistema. A Seção 2 apresenta os conceitos e trabalhos relacionados; a Seção 3 descreve o artefato, o dataset e o experimento; a Seção 4 relata as provas técnicas e o piloto de propostas, discute seus limites e define a apresentação dos resultados finais; e a Seção 5 sintetiza as conclusões permitidas nesta etapa.
